\documentclass[aspectratio=169,10pt]{beamer}

\usetheme{Madrid}
\usecolortheme{default}

\usepackage[T1]{fontenc}
\usepackage[utf8]{inputenc}
\usepackage{lmodern}
\usepackage{booktabs}
\usepackage{array}
\usepackage{tabularx}
\usepackage{amsmath}
\usepackage{amsfonts}
\usepackage{mathtools}
\usepackage{listings}
\usepackage{xcolor}

\definecolor{codebg}{RGB}{248,248,248}
\definecolor{darkblue}{RGB}{25,60,110}

\lstdefinestyle{plaincode}{
  basicstyle=\ttfamily\scriptsize,
  backgroundcolor=\color{codebg},
  columns=fullflexible,
  keepspaces=true,
  breaklines=true,
  frame=single,
  framerule=0.2pt,
  rulecolor=\color{black!20},
  xleftmargin=0.5em,
  xrightmargin=0.5em
}
\lstset{style=plaincode}

\newcommand{\code}[1]{\texttt{\detokenize{#1}}}
\newcommand{\smallcite}[1]{{\scriptsize #1}}

\title[SAQ Follow-Up Progress]{SAQ Follow-Up Research Progress}
\subtitle{From Local Plan Corrections To Graph-Traversal Compatibility}
\author{Kaibin Lu}
\institute{PhD Research Meeting}
\date{2026-07-09; research-validity revision 2026-07-10}

\begin{document}

% ----------------------------------------------------------------------
\begin{frame}
  \titlepage
\end{frame}

% ----------------------------------------------------------------------
\begin{frame}[fragile]{Meeting Goal}
\small
This meeting is not about presenting one finished method.

\vspace{0.6em}
The goal is to explain, from a PhD research perspective:
\begin{enumerate}
  \item What SAQ already contributes.
  \item Why a follow-up is difficult.
  \item What directions we tried.
  \item Which directions produced useful evidence but should not be the main method.
  \item What the current graph-traversal direction is testing.
  \item What evidence is still missing before this can become a SIGMOD/VLDB/ICDE-level contribution.
\end{enumerate}

\vspace{0.4em}
\begin{block}{Main message}
\small
So far, we have not found an independent method that clearly dominates SAQ.
The useful research value is that these attempts narrow the search space:
small empirical plan tuning is weak, mixed local plans add overhead, static
global cost objectives do not transfer, and SAQ's IVF search path is already
hard to improve with simple query-unaware scheduling.
\end{block}
\end{frame}

% ----------------------------------------------------------------------
\begin{frame}[fragile]{Branch Map}
\scriptsize
\begin{tabularx}{\textwidth}{>{\ttfamily}lX}
\toprule
branch & role \\
\midrule
main & upstream SAQ code \\
saq-correctness-base & clean base with only confirmed correctness fixes \\
saq-boundary-audit & older default-neighborhood fixed-policy experiments \\
saq-structural-followup & mixed shared local-plan experiments \\
saq-global-cost-dp & one-global-plan static segment-cost DP experiments \\
saq-planner-objective-analysis & fac-error planner objective and IVF search-procedure measurements \\
saq-graph-traversal-analysis & current graph-index compatibility direction \\
\bottomrule
\end{tabularx}

\vspace{0.8em}
\begin{block}{Current branch}
\begin{lstlisting}
saq-graph-traversal-analysis
validity plan recorded at: a33f7fe
\end{lstlisting}
\end{block}

\small
Except for \code{main} and \code{saq-correctness-base}, the other branches are
research-attempt branches. The old fixed-policy slides are stylistically useful
but substantively out of date.
\end{frame}

% ----------------------------------------------------------------------
\begin{frame}[fragile]{What We Can Safely Say Today}
\small
\begin{block}{Safe statement}
SAQ is a strong baseline. Most simple query-unaware corrections around its
segment planner or IVF search path are either too empirical, too small, or too
expensive.
\end{block}

\vspace{0.5em}
\begin{itemize}
  \item Fixed-policy local plan corrections: repeated small positives on GIST/CIFAR, but empirical and handcrafted.
  \item Shared local plans: residual-local signal exists, but mixed-plan dispatch loses QPS.
  \item Static global segment-cost DP: SAQ default is not dominated on simple risk-cost frontiers.
  \item Fac-error planner objective: plan can move, but GIST custom plan does not recover default recall before losing the QPS advantage.
  \item Search-procedure measurements: fast-stage pruning and accurate-stage early exit are already effective.
  \item Graph-traversal direction: the old unrotated formula-level proxy
  comparison is provisional. A review-time FHT reproduction may remove or
  reverse the apparent SAQ local-rank advantage.
\end{itemize}

\begin{alertblock}{Current graph-evidence boundary}
No current result supports an SAQ-versus-SymphonyQG conclusion or a
path-dependent graph-traversal claim.
\end{alertblock}
\end{frame}

% ----------------------------------------------------------------------
\begin{frame}[fragile]{Core Problem Setting}
\small
We store a database:
\[
X = \{x_1, x_2, \ldots, x_N\}, \quad x_i \in \mathbb{R}^D.
\]

Given a query:
\[
q \in \mathbb{R}^D,
\]
nearest-neighbor search wants:
\[
\arg\min_i d(q, x_i).
\]

For L2 search:
\[
d(q,x_i)=\|q-x_i\|_2^2.
\]

\vspace{0.5em}
In high-dimensional ANNS, exact float32 distance computation is expensive:
many candidates may be scanned per query, and each candidate has \(D\) dimensions.
Vector quantization stores compressed codes and estimates distances from these codes.
\end{frame}

% ----------------------------------------------------------------------
\begin{frame}[fragile]{SAQ In One Slide}
\small
SAQ stands for Segmented Adjusted Quantization.

\vspace{0.5em}
High-level flow:
\begin{enumerate}
  \item Rotate data with PCA.
  \item Split PCA dimensions into contiguous segments.
  \item Assign more bits to high-variance segments and fewer bits to low-variance segments.
  \item Use CAQ-style code adjustment inside each segment.
  \item Use staged distance estimation during search.
\end{enumerate}

\vspace{0.6em}
\begin{block}{Important code paths}
\begin{lstlisting}
saqlib/quantization/saq_data.hpp
  dynamic_programming(...)

saqlib/quantization/saq_estimator.hpp
  varsEstDist(...)
  compFastDist(...)
  compAccurateDist(...)

saqlib/quantization/saq_searcher.hpp
  progressive block scan and refinement
\end{lstlisting}
\end{block}
\end{frame}

% ----------------------------------------------------------------------
\begin{frame}[fragile]{Segment Plan Notation}
\small
We write a plan as:
\begin{lstlisting}
dim:bits, dim:bits, ...
\end{lstlisting}

Running example, GIST full / K4096 / B=4:
\begin{lstlisting}
64:11,192:6,320:4,256:2,128:0
\end{lstlisting}

\vspace{0.5em}
\scriptsize
\begin{tabular}{rrrr}
\toprule
segment & PCA dimension range & dimensions & bits/dim \\
\midrule
1 & 0--63 & 64 & 11 \\
2 & 64--255 & 192 & 6 \\
3 & 256--575 & 320 & 4 \\
4 & 576--831 & 256 & 2 \\
5 & 832--959 & 128 & 0 \\
\bottomrule
\end{tabular}

\vspace{0.6em}
\small
The final 0-bit segment is a dropped low-variance tail.
\end{frame}

% ----------------------------------------------------------------------
\begin{frame}[fragile]{Variable Glossary}
\scriptsize
\begin{tabularx}{\textwidth}{lXl}
\toprule
symbol & meaning & typical/default value \\
\midrule
\(N\) & number of base vectors & dataset-dependent \\
\(D\) & original dimension & GIST 960, CIFAR 512, DEEP 256, audio 192 \\
\(D_p\) & padded dimension & rounded to multiple of 64 \\
\(g\) & SAQ segment granularity & 64 dimensions \\
\(n\) & number of 64-d blocks, \(D_p/g\) & GIST 15 \\
\(B\) & average bits per dimension & usually 3, 4, or 5 \\
\(K\) & IVF cluster count & GIST full 4096, many samples 512 \\
\(S\) & number of SAQ segments & plan-dependent \\
\(b_s\) & bitwidth of segment \(s\) & 0 to 13 \\
\(P_0\) & SAQ default segment plan & DP output \\
\(P\) & candidate/custom segment plan & direction-dependent \\
\(k\) & retrieval cutoff & R@100 or R@10 \\
\(\text{nprobe}\) & number of IVF clusters scanned & setting-dependent \\
\bottomrule
\end{tabularx}

\vspace{0.8em}
\begin{block}{Code constants}
\begin{lstlisting}
kDimPaddingSize = 64
KMaxQuantizeBits = 13
KFastScanSize = 32
positive-segment factor overhead = 64 bits
\end{lstlisting}
\end{block}
\end{frame}

% ----------------------------------------------------------------------
\begin{frame}[fragile]{SAQ Default Planner Formula}
\small
SAQ's default dynamic program minimizes a variance-risk proxy.

\vspace{0.5em}
For a positive-bit segment \([a,b)\) with bitwidth \(r\):
\[
\operatorname{risk}([a,b),r)
= \frac{\sum_{i=a}^{b-1}\operatorname{var}_i}{2^r}.
\]

For a zero-bit tail:
\[
\operatorname{risk}([a,D),0)
= \sum_{i=a}^{D-1}\operatorname{var}_i.
\]

Conceptual DP state:
\[
f[\text{num\_segments}][\text{block\_index}][\text{used\_bits}]
= \text{minimum risk so far}.
\]

\begin{block}{Code}
\begin{lstlisting}
saqlib/quantization/saq_data.hpp
  SaqDataMaker::dynamic_programming(...)
\end{lstlisting}
\end{block}
\end{frame}

% ----------------------------------------------------------------------
\begin{frame}[fragile]{SAQ Planner Complexity}
\small
Let:
\[
n = D_p/64,\quad C = \lfloor B D_p + 64 \rfloor,\quad R = 13,
\]
and let \(S_{\max}\) be the maximum number of segments considered.

\vspace{0.5em}
The implementation enumerates:
\[
\text{segment count} \times \text{current block} \times
\text{used bits} \times \text{next segment length} \times \text{bitwidth}.
\]

Conservative complexity:
\[
\text{time} = O(S_{\max} \cdot n \cdot C \cdot n \cdot R),
\]
\[
\text{memory} = O(S_{\max} \cdot n \cdot C).
\]

For GIST B=4:
\[
D_p=960,\quad n=15,\quad C=4\cdot960+64=3904,\quad R=13.
\]

\vspace{0.4em}
This DP is offline and cheap compared with full index construction, but it only
optimizes the variance-risk proxy.
\end{frame}

% ----------------------------------------------------------------------
\begin{frame}[fragile]{Running Example: Why The Follow-Up Is Hard}
\small
SAQ default for GIST full / K4096 / B=4:
\begin{lstlisting}
P0 = 64:11,192:6,320:4,256:2,128:0
\end{lstlisting}

\vspace{0.5em}
A tempting idea:
\begin{lstlisting}
Maybe a nearby plan can be faster or more accurate.
\end{lstlisting}

\vspace{0.5em}
But any custom plan must answer:
\begin{enumerate}
  \item Does it improve or preserve recall?
  \item Does it improve QPS or memory traffic?
  \item Does it avoid extra index/search overhead?
  \item Does it generalize beyond one dataset?
  \item Is the selection rule query-unaware?
  \item Is it more than parameter tuning around SAQ?
\end{enumerate}
\end{frame}

% ----------------------------------------------------------------------
\begin{frame}[fragile]{Correctness Base: Not A Research Contribution}
\scriptsize
The clean base branch preserves two correctness fixes:
\begin{enumerate}
  \item positive 1-bit segment packing support;
  \item finite valid-lane SIMD block-min default.
\end{enumerate}

Relevant modes:
\begin{lstlisting}
default: -searcher_safe_block_min_mode=2
legacy comparison only: -searcher_safe_block_min_mode=0
\end{lstlisting}

Meaning:
\begin{lstlisting}
0 = legacy native behavior
1 = scalar finite valid-lane block minimum
2 = SIMD finite valid-lane block minimum (default)
\end{lstlisting}

Code paths:
\begin{lstlisting}
src/define_options.h
saqlib/quantization/block_min.hpp
saqlib/quantization/saq_searcher.hpp
unit_test/ut_searcher_correctness.cpp
\end{lstlisting}

Phase 1 verifies valid-lane counts 1--31, positive 1-bit index round trip,
Release/ASAN builds, and full GIST/K4096/B=3 build/load/search. These fixes are
necessary for fair evaluation, but are not research novelty.
\end{frame}

% ----------------------------------------------------------------------
\section{Attempt 1: Fixed Policy}

\begin{frame}[fragile]{Attempt 1: Default-Neighborhood Fixed Policy}
\small
Research question:
\begin{lstlisting}
Can we improve SAQ by searching a small local neighborhood around its default
plan using only base/index artifacts, without representative query workloads?
\end{lstlisting}

Method:
\begin{lstlisting}
SAQ default plan P0
  -> classify default-plan shape
  -> generate local candidate neighborhood N(P0)
  -> score candidates with data-only boundary pairs
  -> promote, reject, or abstain
\end{lstlisting}

Core boundary-pair weight:
\[
w = \exp(-m/\tau),
\quad
m=d_{\mathrm{exact}}(\text{negative})-d_{\mathrm{exact}}(\text{positive}).
\]

\begin{block}{Old branch code}
\begin{lstlisting}
script/generate_default_neighborhood_plans.py
script/score_default_neighborhood_plans.py
script/sweep_data_boundary_pairs.py
script/run_fixed_policy_matrix.py
\end{lstlisting}
\end{block}
\end{frame}

% ----------------------------------------------------------------------
\begin{frame}[fragile]{Fixed Policy: Running Example}
\small
GIST full / K4096 / B=4.

\begin{columns}[T,totalwidth=\textwidth]
\begin{column}{0.48\textwidth}
\begin{block}{Default}
\begin{lstlisting}
64:11,192:6,320:4,256:2,128:0
\end{lstlisting}
\end{block}
\begin{block}{Selected candidate}
\begin{lstlisting}
128:9,320:5,320:3,192:0
\end{lstlisting}
\end{block}
\end{column}
\begin{column}{0.48\textwidth}
\begin{itemize}
  \item head widened: 64 to 128 dimensions
  \item positive segments reduced: 4 to 3
  \item zero tail expanded: 128 to 192 dimensions
\end{itemize}
\end{column}
\end{columns}

\vspace{0.8em}
\begin{block}{Scorer and validation}
\begin{lstlisting}
recall-risk = 0.9071
speed-proxy = 0.7792
role = conservative_eligible

R@100 at nprobe 800: 0.98845 -> 0.98922   delta +0.00077
QPS at nprobe 800:   1013.64 -> 1211.10   ratio 1.1948x
\end{lstlisting}
\end{block}
\end{frame}

% ----------------------------------------------------------------------
\begin{frame}[fragile]{Fixed Policy: Evidence Table}
\scriptsize
\begin{tabularx}{\textwidth}{l l X X}
\toprule
setting & decision & selected/tested plan & measured readout \\
\midrule
GIST K4096 B=3 & promote & \code{64:8,320:5,320:2,256:0} & +0.00159 R@100, 1.1119x QPS at np800 \\
GIST K4096 B=4 & promote & \code{128:9,320:5,320:3,192:0} & +0.00077 R@100, 1.1948x QPS at np800 \\
GIST K4096 B=5 & promote & \code{128:9,128:7,320:5,320:3,64:0} & +0.00028 R@100, 1.0793x QPS at np800 \\
CIFAR B=3 & promote & \code{128:6,64:4,192:2,128:0} & +0.0004 R@10, 1.0761x QPS at np200 \\
CIFAR B=4 & promote & \code{128:7,256:4,128:0} & +0.0004 R@10, 1.0745x QPS at np200 \\
CIFAR B=5 & promote & \code{128:8,64:6,256:4,64:0} & +0.0008 R@10, 1.0609x QPS at np200 \\
DEEP B=4 & reject & \code{128:4,128:3} & +QPS, but -0.02757 R@100 at np200 \\
DEEP B=5 & reject & \code{128:5,128:4} & +QPS, but -0.01429 R@100 at np200 \\
audio B=4 & abstain & none & single-uniform default \\
word2vec B=4 & abstain & none & single-uniform default \\
\bottomrule
\end{tabularx}
\end{frame}

% ----------------------------------------------------------------------
\begin{frame}[fragile]{Fixed Policy: Complexity And Overhead}
\small
Let:
\[
C_{\mathrm{cand}}=\text{number of candidate plans},\quad
P=\text{number of boundary pairs},\quad S=\text{segments}.
\]

Approximate scoring complexity:
\[
O(C_{\mathrm{cand}}\cdot P \cdot S)
\]
after feature construction.

\vspace{0.5em}
\scriptsize
\begin{tabular}{lrrrr}
\toprule
run & candidates & pairs & planner runtime & qps curve geomean \\
\midrule
GIST K4096 B=3 & 2 & 14740 & 145.265 s & 1.1077 \\
GIST K4096 B=4 & 5 & 14740 & 150.580 s & 1.1511 \\
GIST K4096 B=5 & 5 & 14740 & 149.590 s & 1.0668 \\
CIFAR B=3 & 4 & 1068 & 8.562 s & 1.0528 \\
DEEP B=4 & 2 & 1904 & 5.804 s & 1.0802, rejected \\
\bottomrule
\end{tabular}

\vspace{0.8em}
\small
Conclusion: the scorer cost is large relative to the small metric gains.
This weakens novelty and practical value.
\end{frame}

% ----------------------------------------------------------------------
\section{Attempt 2: Shared Local Plans}

\begin{frame}[fragile]{Attempt 2: Mixed Shared Local Plans}
\small
Research question:
\begin{lstlisting}
Can IVF clusters with different residual profiles share a small number of
local SAQ plans, improving local residual quantization without learning from
queries?
\end{lstlisting}

Core assignment rule:
\[
\operatorname{assign}(c)=\arg\min_P \operatorname{cost}(\operatorname{residual\_profile}_c,P).
\]

Variables:
\begin{itemize}
  \item \(c\): IVF cluster id.
  \item \(P\): one shared plan from a small family.
  \item \(M\): number of shared plans, tested \(M\in\{2,4,8\}\).
\end{itemize}

\begin{block}{Old branch code}
\begin{lstlisting}
script/cluster_residual_feasibility_matrix.py
shared-plan materialization code in saqlib/ and src/
\end{lstlisting}
\end{block}
\end{frame}

% ----------------------------------------------------------------------
\begin{frame}[fragile]{Shared Local Plans: Offline Signal}
\scriptsize
\begin{tabularx}{\textwidth}{l r X r r r r l}
\toprule
case & B & global plan & oracle & M & assigned & retention & decision \\
\midrule
GIST K4096 & 3 & \code{64:9,192:5,320:3,192:1,192:0} & 0.9581 & 4 & 0.9627 & 0.889 & continue \\
GIST K4096 & 4 & \code{64:11,192:6,320:4,256:2,128:0} & 0.9420 & 4 & 0.9435 & 0.974 & continue \\
GIST K4096 & 5 & \code{64:11,192:7,320:5,320:3,64:0} & 0.9606 & 4 & 0.9640 & 0.915 & continue \\
CIFAR B=3 & 3 & \code{64:8,128:4,192:2,128:0} & 0.9641 & 2 & 0.9654 & 0.966 & continue \\
DEEP/audio/word2vec & 3--5 & compact or uniform & near 1.000 & 2 & 1.000 & 0.000 & abstain \\
\bottomrule
\end{tabularx}

\vspace{0.8em}
\small
Interpretation:
\begin{lstlisting}
GIST and some CIFAR cases show local residual structure.
DEEP, audio, and word2vec do not show useful local-plan opportunity.
\end{lstlisting}
\end{frame}

% ----------------------------------------------------------------------
\begin{frame}[fragile]{Shared Local Plans: End-To-End Result}
\scriptsize
GIST full / K4096 / B=4, safe search, R@100, nprobe 200.

\vspace{0.5em}
\begin{tabular}{lrrrr}
\toprule
index & R@100 & QPS & QPS ratio & index size \\
\midrule
SAQ default & 0.94469 & 2677.16 & 1.0000 & 570M \\
shared local M4 & 0.94548 & 2434.73 & 0.9094 & 573M \\
cost-aware M4 & 0.94485 & 2398.13 & 0.8958 & 567M \\
\bottomrule
\end{tabular}

\vspace{1em}
Runtime decomposition:

\begin{tabular}{lrrrr}
\toprule
index & fast bit ratio & accurate bit ratio & accurate segment eval ratio & distinct plans/query \\
\midrule
shared local M4 & 0.9741 & 0.9097 & 0.9853 & 2.907 \\
cost-aware M4 & 0.9547 & 0.8809 & 0.8989 & 3.336 \\
\bottomrule
\end{tabular}

\vspace{1em}
\small
Measured work decreases, but QPS also decreases. The likely cause is mixed-plan
overhead: multiple query-side estimator states, lookup-table setup paths, and
worse cache/layout locality.
\end{frame}

% ----------------------------------------------------------------------
\begin{frame}[fragile]{Shared Local Plans: Complexity}
\small
Let:
\[
K=\text{IVF clusters},\quad M=\text{shared plans},\quad
n=D_p/64,\quad C=\text{bit budget},\quad R=13.
\]

Local per-cluster DP can cost:
\[
O(K\cdot S_{\max}\cdot n\cdot C\cdot n\cdot R).
\]

Shared assignment after plan generation:
\[
O(K\cdot M\cdot \operatorname{cost\_eval}).
\]

Search-time architecture:
\begin{lstlisting}
default SAQ: one plan, one query-side estimator state
shared plans: up to M plan-specific estimator states per query
\end{lstlisting}

Observed on GIST B=4:
\begin{lstlisting}
2.907 to 3.336 distinct plan states per query
\end{lstlisting}

Conclusion: this is useful limitation evidence, but not a main method.
\end{frame}

% ----------------------------------------------------------------------
\section{Attempt 3: Global Segment-Cost DP}

\begin{frame}[fragile]{Attempt 3: Single-Global Segment-Cost DP}
\small
Research question:
\begin{lstlisting}
Can we keep one global SAQ plan but add search-time cost awareness to the
planner objective?
\end{lstlisting}

To avoid arbitrary weights, we did not start with:
\[
\operatorname{risk} + \lambda\cdot\operatorname{cost}.
\]

Instead, we built Pareto frontiers for separate static cost terms:
\begin{itemize}
  \item positive dimensions;
  \item total code-bit volume;
  \item accurate extra bit volume;
  \item positive segment count;
  \item segment-factor bits;
  \item zero-tail risk.
\end{itemize}

\begin{block}{Old branch code}
\begin{lstlisting}
script/global_segment_cost_frontier.py
\end{lstlisting}
\end{block}
\end{frame}

% ----------------------------------------------------------------------
\begin{frame}[fragile]{Global Cost DP: Offline Frontier Result}
\scriptsize
B=4 snapshot:

\vspace{0.4em}
\begin{tabularx}{\textwidth}{lXrrrr}
\toprule
dataset & SAQ default plan & pos dim & pos segs & code bits & closest lower-code risk \\
\midrule
GIST K4096 & \code{64:11,192:6,320:4,256:2,128:0} & 832 & 4 & 3648 & 1.0100 \\
GIST sample100k & same & 832 & 4 & 3648 & 1.0099 \\
CIFAR60K & \code{64:9,192:5,128:3,128:0} & 384 & 3 & 1920 & 1.1104 \\
DEEP sample100k & \code{64:6,192:3} & 256 & 2 & 960 & 1.0164 \\
audio K4096 & \code{192:4} & 192 & 1 & 768 & 2.0000 \\
word2vec sample100k & \code{320:4} & 320 & 1 & 1280 & 1.0319 \\
\bottomrule
\end{tabularx}

\vspace{1em}
\small
Main result:
\begin{lstlisting}
No lower-cost plan dominated the SAQ default at no worse variance-risk.
\end{lstlisting}
\end{frame}

% ----------------------------------------------------------------------
\begin{frame}[fragile]{Global Cost DP: End-To-End Check}
\scriptsize
GIST full / K4096 / B=4, safe search, R@100.

\vspace{0.4em}
\begin{tabularx}{\textwidth}{lX}
\toprule
label & plan \\
\midrule
default & \code{64:11,192:6,320:4,256:2,128:0} \\
code-volume & \code{64:10,128:7,192:5,192:3,256:2,128:0} \\
positive-dim & \code{64:10,128:7,192:5,384:3,192:0} \\
\bottomrule
\end{tabularx}

\vspace{0.8em}
\begin{tabular}{r l r r r}
\toprule
nprobe & label & R@100 & QPS & QPS ratio \\
\midrule
100 & default & 0.86604 & 4275.79 & 1.0000 \\
100 & code-volume & 0.86630 & 4031.72 & 0.9429 \\
100 & positive-dim & 0.86639 & 4212.68 & 0.9852 \\
200 & default & 0.94469 & 2732.35 & 1.0000 \\
200 & code-volume & 0.94506 & 2503.10 & 0.9161 \\
200 & positive-dim & 0.94537 & 2673.65 & 0.9785 \\
400 & default & 0.97999 & 1642.83 & 1.0000 \\
400 & code-volume & 0.98092 & 1499.40 & 0.9127 \\
400 & positive-dim & 0.98139 & 1620.92 & 0.9867 \\
\bottomrule
\end{tabular}

\vspace{0.6em}
\small
Slight recall gain, no QPS gain. Static cost proxies do not transfer reliably.
\end{frame}

% ----------------------------------------------------------------------
\begin{frame}[fragile]{Global Cost DP: Complexity}
\small
Offline frontier generation:
\begin{lstlisting}
enumerate SAQ DP candidate plans
evaluate variance risk and each static cost metric
filter metric-specific Pareto frontiers
\end{lstlisting}

Approximate complexity:
\[
\text{DP candidate enumeration}=O(S_{\max} n C n R),
\]
\[
\text{frontier filtering}=O(P\log P)\quad\text{per cost metric}.
\]

Where:
\[
P = \text{number of retained candidate states/plans}.
\]

\vspace{0.8em}
Why the direction stopped:
\begin{lstlisting}
The complexity is acceptable.
The issue is objective validity: static code-volume or segment-count terms did
not predict real QPS benefit.
\end{lstlisting}
\end{frame}

% ----------------------------------------------------------------------
\section{Attempt 4: Fac-Error Planner Objective}

\begin{frame}[fragile]{Attempt 4: Fac-Error Planner Objective}
\small
Research question:
\begin{lstlisting}
Does SAQ's variance-risk proxy mismatch measured CAQ/SAQ segment error?
\end{lstlisting}

Measured quantities:
\begin{itemize}
  \item variance proxy: \(\sum \operatorname{variance}/2^{\text{bits}}\);
  \item raw SSE;
  \item scale-aligned SSE;
  \item direction loss: \(1-\cos(o,o_a)^2\);
  \item CAQ fac-error.
\end{itemize}

\begin{block}{Old branch code}
\begin{lstlisting}
src/measure_planner_proxy.cpp
script/summarize_planner_proxy.py
script/falsify_fac_error_dp.py
\end{lstlisting}
\end{block}

Candidate enumeration:
\begin{lstlisting}
all contiguous 64-dimensional block intervals
all bitwidths from 0 to 13
deterministic 1024-row or 2048-row residual sample
\end{lstlisting}
\end{frame}

% ----------------------------------------------------------------------
\begin{frame}[fragile]{Planner Proxy Measurement Result}
\scriptsize
SAQ's variance proxy strongly ranks energy-weighted residual error.

\vspace{0.4em}
\begin{tabular}{lrrr}
\toprule
dataset & aligned-SSE Spearman & within-bit Spearman & discordant pairs \\
\midrule
audio K4096 & 0.9113 & 1.0000 & 0.0000 \\
CIFAR60K & 0.9624 & 0.9979 & 0.0080 \\
DEEP sample100k & 0.9651 & 1.0000 & 0.0000 \\
GIST sample100k & 0.9684 & 0.9981 & 0.0123 \\
word2vec sample100k & 0.9917 & 0.9452 & 0.0782 \\
\bottomrule
\end{tabular}

\vspace{1em}
Direction-only mismatch:

\begin{tabular}{lrr}
\toprule
dataset & direction-loss Spearman & fac-error Spearman \\
\midrule
audio K4096 & 0.4769 & 0.7143 \\
CIFAR60K & 0.5516 & 0.9505 \\
DEEP sample100k & 0.5189 & 0.8294 \\
GIST sample100k & 0.5234 & 0.9753 \\
word2vec sample100k & 0.7412 & 0.9838 \\
\bottomrule
\end{tabular}

\vspace{0.8em}
\small
Interpretation: the obvious measured-SSE replacement objective is not promising.
The visible gap is direction quality, but its system relevance is unclear.
\end{frame}

% ----------------------------------------------------------------------
\begin{frame}[fragile]{Fac-Error DP Falsification}
\scriptsize
We replaced SAQ variance cost with measured fac-error cost in the same DP.

\vspace{0.4em}
\begin{tabularx}{\textwidth}{lXXc}
\toprule
dataset & variance plan & fac-error plan & changed? \\
\midrule
audio K4096 & \code{192x4} & \code{192x4} & no \\
CIFAR60K & \code{64x9_192x5_128x3_128x0} & \code{128x8_320x3_64x0} & yes \\
DEEP sample100k & \code{64x6_192x3} & \code{64x6_192x3} & no \\
GIST sample100k & \code{64x11_192x6_320x4_256x2_128x0} & \code{192x9_512x4_256x0} & yes \\
word2vec sample100k & \code{320x4} & \code{320x4} & no \\
\bottomrule
\end{tabularx}

\vspace{1em}
\small
This avoided the trivial negative outcome that the objective always reproduces
SAQ. But end-to-end GIST recall-matched evaluation failed.
\end{frame}

% ----------------------------------------------------------------------
\begin{frame}[fragile]{Fac-Error GIST Recall-Matched Check}
\scriptsize
GIST sample100k / K512 / B=4 / R@100 / 24 threads.

\vspace{0.4em}
Default target:
\begin{lstlisting}
SAQ variance, nprobe=200:
R@100 = 0.99132
QPS = 9215.208
\end{lstlisting}

Fac-error custom plan:
\begin{lstlisting}
192x9_512x4_256x0
\end{lstlisting}

\begin{tabular}{lrrr}
\toprule
plan / nprobe & R@100 & QPS & QPS vs default \\
\midrule
SAQ variance / 200 & 0.99132 & 9215.208 & 1.000x \\
fac-error / 200 & 0.99059 & 12334.094 & 1.338x \\
fac-error / 240 & 0.99084 & 10957.260 & 1.189x \\
fac-error / 280 & 0.99091 & 9875.544 & 1.072x \\
fac-error / 300 & 0.99091 & 9422.062 & 1.022x \\
fac-error / 320 & 0.99091 & 8956.573 & 0.972x \\
\bottomrule
\end{tabular}

\vspace{0.6em}
\small
The custom plan does not recover default recall before crossing below default QPS.
\end{frame}

% ----------------------------------------------------------------------
\begin{frame}[fragile]{Fac-Error Direction: Complexity}
\small
Let:
\[
R_s=\text{sampled residual vectors or pairs},\quad
n=D_p/64,\quad
B_w=\text{tested bitwidth count}.
\]

Number of contiguous intervals:
\[
I=\frac{n(n+1)}{2}.
\]

Measurement complexity:
\[
O(R_s\cdot B_w\cdot \sum\text{interval lengths})
\]
or as an upper view:
\[
O(R_s\cdot B_w\cdot n^3\cdot64).
\]

DP after measurement:
\[
O(S_{\max}nCnR).
\]

Why it stopped:
\begin{lstlisting}
Even when measurement moves the plan, the first changed GIST plan produces a
speed/recall tradeoff, not a recall-matched improvement.
\end{lstlisting}
\end{frame}

% ----------------------------------------------------------------------
\section{Attempt 5: Search Procedure}

\begin{frame}[fragile]{Attempt 5: Search-Procedure Analysis}
\small
Research question:
\begin{lstlisting}
Can SAQ's current IVF search path be improved by a query-unaware estimator
scheduling or pruning rule?
\end{lstlisting}

Search path quantities:
\begin{itemize}
  \item variance stage: cheap lower estimate per block;
  \item fast stage: 1-bit estimate over positive segments;
  \item accurate stage: full-bit segment refinement;
  \item result pool: top-k candidate maintenance.
\end{itemize}

\begin{block}{Old branch instrumentation}
\begin{lstlisting}
runtime profile counters in test_qps
profile_segment_order
profile_variance_bound
\end{lstlisting}
\end{block}

Key code:
\begin{lstlisting}
saqlib/quantization/saq_searcher.hpp
saqlib/quantization/saq_estimator.hpp
src/test_qps.cpp
\end{lstlisting}
\end{frame}

% ----------------------------------------------------------------------
\begin{frame}[fragile]{Runtime Profile Evidence}
\scriptsize
GIST sample100k / K512 / B=4 / default plan:
\begin{lstlisting}
64x11_192x6_320x4_256x2_128x0
\end{lstlisting}

\begin{tabular}{rrrrr}
\toprule
nprobe & R@100 & QPS & scanned cand/query & accurate cand/query \\
\midrule
160 & 0.99036 & 10119.473 & 37339.7 & 2696.3 \\
200 & 0.99132 & 9192.221 & 45810.0 & 2767.4 \\
240 & 0.99153 & 8183.555 & 54044.2 & 2805.3 \\
\bottomrule
\end{tabular}

\vspace{0.8em}
\begin{tabular}{rrrrr}
\toprule
nprobe & variance-pruned & fast-pruned & accurate/scanned & early exits/accurate \\
\midrule
160 & 0.14\% & 45.29\% & 7.22\% & 83.91\% \\
200 & 0.27\% & 52.02\% & 6.04\% & 84.32\% \\
240 & 0.45\% & 57.60\% & 5.19\% & 84.52\% \\
\bottomrule
\end{tabular}

\vspace{0.8em}
\small
Fast-stage pruning and accurate-stage early exit already do useful work.
Variance pruning is almost inactive.
\end{frame}

% ----------------------------------------------------------------------
\begin{frame}[fragile]{Segment-Order Counterfactual}
\scriptsize
Question:
\begin{lstlisting}
Is SAQ's default PCA/bit segment order suboptimal for progressive pruning?
\end{lstlisting}

Tested orders:
\begin{lstlisting}
pca, reverse, bit_desc, dim_desc, cost_asc, risk_per_cost_desc
\end{lstlisting}

At nprobe 200, deltas relative to default \code{pca}:

\begin{tabular}{lrrrr}
\toprule
order & recall delta & fast-call delta & accurate-call delta & fast-bit delta \\
\midrule
reverse & -0.00092 & +47.41\% & +167.81\% & +69.06\% \\
bit\_desc & 0.00000 & 0.00\% & 0.00\% & 0.00\% \\
dim\_desc & -0.00091 & +47.39\% & +148.07\% & +69.04\% \\
cost\_asc & +0.00002 & +1.34\% & +0.48\% & +1.78\% \\
risk\_per\_cost\_desc & 0.00000 & 0.00\% & 0.00\% & 0.00\% \\
\bottomrule
\end{tabular}

\vspace{0.8em}
\small
Default PCA/bit order is already tied for lowest work among simple
hyperparameter-free alternatives.
\end{frame}

% ----------------------------------------------------------------------
\begin{frame}[fragile]{Variance-Bound Inactivity}
\scriptsize
Final search-procedure measurement:
\begin{lstlisting}
Why does the variance stage prune less than 0.5% of blocks?
\end{lstlisting}

\begin{tabular}{rrrrr}
\toprule
nprobe & blocks & finite-boundary blocks & variance-pruned & variance-prune rate \\
\midrule
160 & 1,245,446 & 1,241,443 & 1,750 & 0.1405\% \\
200 & 1,530,007 & 1,526,004 & 4,104 & 0.2682\% \\
240 & 1,807,095 & 1,803,092 & 8,173 & 0.4523\% \\
\bottomrule
\end{tabular}

\vspace{0.8em}
Gap to the top-k boundary:

\begin{tabular}{rrr}
\toprule
nprobe & median relative gap & 90th-percentile relative gap \\
\midrule
160 & 0.7467 & 0.8635 \\
200 & 0.7427 & 0.8585 \\
240 & 0.7386 & 0.8537 \\
\bottomrule
\end{tabular}

\vspace{0.8em}
\small
Useful pruning would require removing a large fraction of conservative slack,
which has no safe derivation in the current evidence.
\end{frame}

% ----------------------------------------------------------------------
\begin{frame}[fragile]{Search-Procedure Complexity}
\small
Let:
\[
Q=\text{queries},\quad
B_c=\text{32-lane blocks scanned/query},\quad
S_{\mathrm{pos}}=\text{positive-bit segments},
\]
\[
R_{\mathrm{acc}}=\text{candidates entering accurate refinement},\quad
S_{\mathrm{acc}}=\text{average accurate segments per refined candidate}.
\]

Simplified per-query work:
\[
\text{variance stage}=O(B_c),
\]
\[
\text{fast stage}=O(B_c\cdot\text{segments evaluated before block pruning}),
\]
\[
\text{accurate stage}=O(R_{\mathrm{acc}}\cdot S_{\mathrm{acc}}).
\]

Measured GIST B=4 at nprobe 200:
\begin{lstlisting}
B_c ~= 1530 blocks/query
fast segment calls ~= 4135/query
R_acc ~= 2767 candidates/query
S_acc ~= 1.71 segments/candidate
\end{lstlisting}
\end{frame}

% ----------------------------------------------------------------------
\section{Current Direction: Graph Traversal}

\begin{frame}[fragile]{Pivot: Why Graph Traversal Is Different}
\small
IVF search:
\begin{lstlisting}
centroid probing determines candidate clusters
SAQ estimates mostly filter and rerank vectors inside selected clusters
\end{lstlisting}

Graph ANNS:
\begin{lstlisting}
approximate distance can determine which node is expanded next
\end{lstlisting}

That changes the role of quantization error:
\begin{lstlisting}
IVF error affects scoring among reached candidates.
Graph error can affect the path through the graph.
\end{lstlisting}

\begin{block}{Current research question}
Can SAQ-style progressive compressed distance estimation stabilize graph
frontier decisions, and does it provide a refinement/work tradeoff beyond
RaBitQ/SymphonyQG-style graph quantization?
\end{block}
\end{frame}

% ----------------------------------------------------------------------
\begin{frame}[fragile]{Related Work Gate For Graph Direction}
\small
Important related work:
\begin{itemize}
  \item SymphonyQG;
  \item NGT-QG;
  \item Routing-Guided Learned Product Quantization.
\end{itemize}

What SymphonyQG already covers:
\begin{itemize}
  \item graph-side neighbor code layout;
  \item FastScan distance estimates during graph traversal;
  \item RaBitQ-style graph quantization;
  \item current-vertex residual normalization;
  \item multiple estimates to reduce missed nearest neighbors.
\end{itemize}

Unsafe claims:
\begin{lstlisting}
first quantization + graph index integration
first estimated-distance graph traversal
first graph-specific RaBitQ normalization
\end{lstlisting}

Narrow possible SAQ-specific claim:
\begin{lstlisting}
SAQ's segmented progressive estimator may offer a better frontier refinement
axis than a single-stage RaBitQ/FastScan-style graph estimator.
\end{lstlisting}
\end{frame}

% ----------------------------------------------------------------------
\begin{frame}[fragile]{Graph Replay Measurement Design}
\scriptsize
Diagnostic harness:
\begin{lstlisting}
build fixed exact-kNN adjacency on a deterministic subset
for query q and current graph node u:
  compare exact and compressed estimates over N(u)
\end{lstlisting}

Metric:
\[
v_{\mathrm{exact}}=\arg\min_{v\in N(u)}d_{\mathrm{exact}}(q,v),
\]
\[
\operatorname{rank}_{\mathrm{est}}(v_{\mathrm{exact}})
=\text{rank of }v_{\mathrm{exact}}\text{ under estimator distance}.
\]

Reported:
\begin{itemize}
  \item top1 disagreement rate;
  \item mean and p90 rank of exact-best neighbor;
  \item top-L containment curve;
  \item average distinct IVF clusters per expansion event.
\end{itemize}

Complexity:
\[
\text{adjacency build}=O(M^2D),\quad
\text{replay scoring}=O(E\cdot \text{degree}\cdot\text{estimator\_cost}).
\]

\begin{alertblock}{Current implementation boundary}
Each query-root neighbor set is scored independently. There is no frontier
heap, visited set, path-dependent beam traversal, or SymphonyQG
multiple-estimate behavior. Rows named \code{saq_prefix_accp} accurately score
the first (p) segments for every candidate; they are not top-(r) candidate
refinement.
\end{alertblock}
\end{frame}

% ----------------------------------------------------------------------
\begin{frame}[fragile]{Historical Weak-Proxy GIST Diagnostic}
\scriptsize
GIST sample50k / K512 / B=4:
\begin{lstlisting}
plan = 64:11,192:6,320:4,256:2,128:0
subset = 4096
degree = 32
queries = 100
roots per query = 8
events = 800
\end{lstlisting}

\begin{tabular}{lrrrrr}
\toprule
estimator & code bits only & top1 disag. & mean rank & p90 rank & top8 contain. \\
\midrule
global rabitq-style proxy & 960 & 0.90625 & 12.4812 & 27 & 0.44625 \\
saq\_fast & 832 & 0.46125 & 2.2575 & 5 & 0.98125 \\
saq\_prefix\_acc1 & 1472 & 0.13750 & 1.1725 & 2 & 1.00000 \\
saq\_prefix\_acc2 & 2432 & 0.02750 & 1.0300 & 1 & 1.00000 \\
saq\_full & 3648 & 0.00625 & 1.00625 & 1 & 1.00000 \\
\bottomrule
\end{tabular}

\vspace{0.6em}
\small
SAQ staged estimates ordered these independent local neighbor sets better than
the weak global proxy. This was not a faithful SymphonyQG comparison and does
not measure path-dependent graph traversal. The code-bit column omits factors,
padding, metadata, estimator preparation, and runtime.
\end{frame}

% ----------------------------------------------------------------------
\begin{frame}[fragile]{SymphonyQG Sanity Baseline}
\scriptsize
Local official SymphonyQG source:
\begin{lstlisting}
/rwproject/kdd-db/kluaq/SymphonyQG
commit 6124ddb
\end{lstlisting}

Small GIST baseline:
\begin{lstlisting}
N = 10,000 base vectors
Q = 100 queries
D = 960
top-k = 10
degree = 32
build EF = 100
build iterations = 2
\end{lstlisting}

\begin{tabular}{rrr}
\toprule
search EF & Recall@10 & QPS \\
\midrule
20 & 0.891 & 37173.659 \\
40 & 0.960 & 29269.393 \\
80 & 0.983 & 20404.281 \\
120 & 0.994 & 16214.257 \\
200 & 0.997 & 10358.608 \\
\bottomrule
\end{tabular}

\vspace{0.8em}
\small
This is not directly comparable to SAQ replay, but it confirms the official
implementation is usable for source-level baseline extraction.
\end{frame}

% ----------------------------------------------------------------------
\begin{frame}[fragile]{SymphonyQG Estimator Formula}
\small
For current graph vertex \(c\), outgoing neighbor \(x\), and query \(q\), the
official path first pads to a power-of-two dimension and applies a shared
random-sign FHT. The following quantities live in that transformed space:
\[
r=x-c,\qquad s_i=
\begin{cases}
+1,& r_i>0,\\
-1,& r_i\le 0.
\end{cases}
\]

\[
\operatorname{fac\_norm}=1/\sqrt{D},
\quad
\operatorname{fac\_x0}=\langle r,s\cdot\operatorname{fac\_norm}\rangle/\|r\|,
\]
\[
\operatorname{fac\_x1}=\langle c,s\cdot\operatorname{fac\_norm}\rangle,
\quad
x\_x0=\|r\|/\operatorname{fac\_x0}.
\]

\[
\operatorname{triple\_x}=\|r\|^2+2\cdot x\_x0\cdot\operatorname{fac\_x1},
\]
\[
\operatorname{factor\_dq}=-2\cdot x\_x0\cdot\operatorname{fac\_norm},
\quad
\operatorname{factor\_vq}=\operatorname{factor\_dq}\sum_i s_i.
\]
\end{frame}

% ----------------------------------------------------------------------
\begin{frame}[fragile]{SymphonyQG Query Quantization And Distance}
\small
Query is scalar-quantized with \(QG\_BQUERY=6\):
\[
\operatorname{width}=\frac{\max(q)-\min(q)}{2^6-1},
\]
\[
q\_code_i=\operatorname{round}\left(\frac{q_i-\min(q)}{\operatorname{width}}+0.5\right).
\]

Signed dot product:
\[
\operatorname{dot\_code}=\sum_i s_i q\_code_i.
\]

Extracted distance:
\[
d(q,x\mid c)=\|q-c\|^2+\operatorname{triple\_x}
\]
\[
\quad+\operatorname{factor\_dq}\cdot\operatorname{width}\cdot\operatorname{dot\_code}
\]
\[
\quad+\operatorname{factor\_vq}\cdot\min(q).
\]

Source:
\begin{lstlisting}
SymphonyQG/symqglib/quantization/rabitq.hpp
SymphonyQG/symqglib/qg/qg_query.hpp
SymphonyQG/symqglib/qg/qg_scanner.hpp
SymphonyQG/symqglib/qg/qg.hpp
\end{lstlisting}
\end{frame}

% ----------------------------------------------------------------------
\begin{frame}[fragile]{Current Incomplete Formula-Level Proxy}
\small
Implemented in current branch:
\begin{lstlisting}
src/profile_graph_frontier.cpp
  SymphonyQGVertexProxy
\end{lstlisting}

What it preserves:
\begin{itemize}
  \item current-vertex residual centering;
  \item 1-bit residual signs;
  \item 6-bit query scalar quantization;
  \item \(\operatorname{triple\_x}\), \(\operatorname{factor\_dq}\), \(\operatorname{factor\_vq}\) formula.
\end{itemize}

What it omits:
\begin{itemize}
  \item FHT rotation;
  \item padded power-of-two dimension;
  \item packed FastScan layout;
  \item SIMD lookup table reproduction;
  \item SymphonyQG-built graph;
  \item multiple estimates for a vertex reached from different parents.
\end{itemize}

Complexity:
\[
\text{formula-level proxy}=O(D)\text{ per edge}.
\]

\begin{alertblock}{Validity limitation}
FHT and padding affect estimator geometry and rank quality; they are not merely
performance optimizations. This proxy is not a source-aligned SymphonyQG
baseline.
\end{alertblock}
\end{frame}

% ----------------------------------------------------------------------
\begin{frame}[fragile]{Provisional Formula-Level Proxy Result}
\scriptsize
GIST sample50k / K512 / B=4 / subset 4096.

\begin{tabular}{lrrrrrr}
\toprule
estimator & code bits only & top1 disag. & mean rank & p90 rank & top4 & top8 \\
\midrule
global rabitq-style proxy & 960 & 0.90625 & 12.4812 & 27 & 0.25625 & 0.44625 \\
symqg\_vertex\_proxy & 960 & 0.87625 & 8.2775 & 19 & 0.38250 & 0.61500 \\
saq\_fast & 832 & 0.46125 & 2.2575 & 5 & 0.87625 & 0.98125 \\
saq\_prefix\_acc1 & 1472 & 0.13750 & 1.1725 & 2 & 0.99750 & 1.00000 \\
saq\_prefix\_acc2 & 2432 & 0.02750 & 1.0300 & 1 & 1.00000 & 1.00000 \\
saq\_full & 3648 & 0.00625 & 1.00625 & 1 & 1.00000 & 1.00000 \\
\bottomrule
\end{tabular}

\vspace{1em}
\small
Interpretation:
\begin{lstlisting}
Within this incomplete unrotated comparison, the vertex-centered proxy improves
over the global proxy and SAQ has better local-neighborhood rank metrics.
\end{lstlisting}

Strict limitation:
\begin{lstlisting}
This table cannot support an SAQ-versus-SymphonyQG claim or an end-to-end graph
claim. Its code-bit counts are not complete storage or runtime work.
\end{lstlisting}
\end{frame}

% ----------------------------------------------------------------------
\begin{frame}[fragile]{Source-Aligned Scalar Baseline Correction}
\scriptsize
A committed scalar implementation now matches pinned SymphonyQG revision
\texttt{6124ddb3...} field by field. The first GIST sample50k / K512 / B=4 /
subset-512 / seed-0 sanity run has 128 events:

\begin{tabular}{lrrrrr}
\toprule
estimator & code bits & top1 disag. & mean rank & p90 rank & top4 \\
\midrule
unrotated symqg proxy & 960 & 0.867188 & 8.85156 & 22 & 0.390625 \\
symqg\_fht\_scalar & 1024 & 0.304688 & 1.54688 & 3 & 0.968750 \\
saq\_fast & 832 & 0.453125 & 1.93750 & 4 & 0.914062 \\
\bottomrule
\end{tabular}

\vspace{0.8em}
\begin{alertblock}{Interpretation}
The aligned scalar baseline is much stronger than the old proxy and is stronger
than saq\_fast in this small setting. The old comparison is invalid as
comparative evidence.
\end{alertblock}

Release and ASAN parity cover query state, edge factors, distances, and tie
ordering. Packed FastScan, multiple estimates, and traversal remain absent.
\end{frame}

% ----------------------------------------------------------------------
\begin{frame}[fragile]{Current Direction: Fuller SymphonyQG Alignment}
\scriptsize
Immediate research question:
\begin{lstlisting}
Does SAQ's segmented progressive estimator retain a graph-frontier
rank-recovery advantage over a closer SymphonyQG/RaBitQ FastScan-style
estimator?
\end{lstlisting}

Required ordered work:
\begin{enumerate}
  \item Phase 1: completed; correctness defaults and focused tests pass.
  \item Phase 2: scalar milestone complete; parity-tested packed path remains.
  \item Phase 3: fixed-seed same-replay evaluation and stop/continue decision.
\end{enumerate}

Phase 2 must reproduce:
\begin{itemize}
  \item FHT rotation;
  \item padded dimension;
  \item 6-bit query quantization;
  \item current-vertex residual neighbor code;
  \item \(\operatorname{triple\_x}/\operatorname{factor\_dq}/\operatorname{factor\_vq}\);
  \item FastScan-equivalent signed dot product.
\end{itemize}

Same-replay comparison:
\begin{lstlisting}
exact_float
symqg_vertex_proxy       historical diagnostic only
symqg_fht_scalar         source-aligned reference
symqg_fht_fastscan       after parity testing
saq_fast
saq_prefix_acc1/2
saq_full
\end{lstlisting}
\end{frame}

% ----------------------------------------------------------------------
\begin{frame}[fragile]{Overall Evidence Summary}
\scriptsize
\begin{tabularx}{\textwidth}{lXX}
\toprule
direction & evidence & status \\
\midrule
fixed-policy default-neighborhood & GIST/CIFAR small positives, DEEP reject, audio/word2vec abstain; scorer overhead 145--151s on GIST & useful but too empirical \\
mixed shared local plans & residual-local signal; GIST B=4 R@100 +0.00079 but QPS ratio 0.9094 & stopped as main method \\
global static segment-cost DP & SAQ default not dominated; near-frontier GIST plans slightly higher recall but slower & stopped as main method \\
fac-error planner objective & objective changes GIST/CIFAR plans; GIST custom plan faster at same nprobe but not recall-matched & stopped as main method \\
search-procedure scheduling & fast pruning and accurate early exit already strong; variance pruning weak but unsafe to tighten empirically & stopped as main method \\
graph traversal compatibility & aligned scalar baseline is stronger than saq\_fast on a small seed-0 sanity run; packed/traversal evidence absent & active only as a source-aligned baseline-validation gate \\
\bottomrule
\end{tabularx}
\end{frame}

% ----------------------------------------------------------------------
\begin{frame}[fragile]{What We Should Not Claim}
\small
Unsafe claims:
\begin{itemize}
  \item We already have a new quantizer.
  \item We already have a method that universally improves SAQ.
  \item The fixed-policy scorer is a paper-level contribution.
  \item The shared local plan direction improves SAQ end to end.
  \item Static segment-cost DP solves SAQ's planner limitation.
  \item Fac-error DP gives a recall-matched speedup.
  \item Simple search-loop scheduling improves SAQ.
  \item The current graph result proves end-to-end graph-search improvement.
  \item A comparative SAQ/SymphonyQG local-order advantage has already been established.
\end{itemize}

\begin{block}{Safer claim}
The negative and partial-positive evidence narrows the credible follow-up
space. The current graph direction tests a structural mismatch between SAQ's
IVF-centered progressive estimator and graph traversal decisions, with
SymphonyQG as the required novelty gate.
\end{block}
\end{frame}

% ----------------------------------------------------------------------
\begin{frame}[fragile]{Proposed Meeting Discussion}
\small
\begin{enumerate}
  \item Is graph-traversal compatibility the right remaining direction, given the negative evidence on local plan tuning and IVF search-loop scheduling?
  \item What level of SymphonyQG baseline alignment is enough before proposing a SAQ-specific graph frontier refinement policy?
  \item If fuller SymphonyQG erases the current advantage, should we pivot to theory of SAQ/CAQ estimator bounds or to transform objectives beyond PCA?
  \item Should the paper story include the negative evidence as a motivation section, or keep it only as internal research reasoning?
\end{enumerate}

\vspace{1em}
\begin{block}{Suggested meeting stance}
Evaluate novelty first, before spending time on end-to-end graph integration.
\end{block}
\end{frame}

% ----------------------------------------------------------------------
\begin{frame}[fragile]{Recommended Next Step}
\small
Do not start full HNSW/DiskANN integration yet.

\vspace{0.6em}
Immediate next step:
\begin{lstlisting}
Phase 2: implement the packed/FastScan-equivalent path and establish ordering
parity with the validated source-aligned scalar estimator.
\end{lstlisting}

Then execute Phase 3:
\begin{lstlisting}
rerun the same local replay over predeclared seeds 0..9
apply the documented stop/continue condition
\end{lstlisting}

Pass condition:
\begin{lstlisting}
At least one SAQ stage is consistently non-dominated in rank quality versus
complete logical work across the predeclared rotations.
\end{lstlisting}

Fail condition:
\begin{lstlisting}
A source-aligned SymphonyQG estimator matches or dominates saq_fast and no SAQ
prefix point gives a stable Pareto improvement after complete work accounting.
\end{lstlisting}
\end{frame}

% ----------------------------------------------------------------------
\appendix
\section{Backup}

\begin{frame}[fragile]{Backup: One-Slide Running Example}
\small
GIST full / K4096 / B=4, SAQ default:
\begin{lstlisting}
64:11,192:6,320:4,256:2,128:0
\end{lstlisting}

Old fixed-policy candidate:
\begin{lstlisting}
128:9,320:5,320:3,192:0
R@100 np800: 0.98845 -> 0.98922
QPS np800:   1013.64 -> 1211.10
\end{lstlisting}

Why not enough:
\begin{lstlisting}
small recall delta
handcrafted local candidate families
145-151s planner/scorer overhead on GIST
not a new quantizer or principled planner
\end{lstlisting}

Current graph running example:
\begin{lstlisting}
GIST sample50k / K512 / B=4 / subset 4096
unrotated symqg_vertex_proxy p90 rank = 19       historical
review-time FHT reproduction p90 rank = 2-3      provisional
saq_fast p90 rank = 5
\end{lstlisting}

Why incomplete:
\begin{lstlisting}
the old comparison omitted source-critical transform steps
the FHT reproduction is not a committed parity-tested runner
neither experiment executes graph traversal
\end{lstlisting}
\end{frame}

% ----------------------------------------------------------------------
\begin{frame}[fragile]{Backup: Code Map}
\scriptsize
Current branch:
\begin{lstlisting}
SAQ planner:
  saqlib/quantization/saq_data.hpp

SAQ estimator:
  saqlib/quantization/saq_estimator.hpp
  saqlib/quantization/caq/caq_estimator.hpp

SAQ search path:
  saqlib/quantization/saq_searcher.hpp
  src/test_qps.cpp

Graph replay:
  src/profile_graph_frontier.cpp

Correctness/evaluation flags:
  src/define_options.h
\end{lstlisting}

Historical branches:
\begin{lstlisting}
fixed policy:
  script/generate_default_neighborhood_plans.py
  script/score_default_neighborhood_plans.py
  script/sweep_data_boundary_pairs.py

shared local plans:
  script/cluster_residual_feasibility_matrix.py

global cost DP:
  script/global_segment_cost_frontier.py

fac-error planner:
  src/measure_planner_proxy.cpp
  script/summarize_planner_proxy.py
  script/falsify_fac_error_dp.py
\end{lstlisting}
\end{frame}

\end{document}
